\input{preamble}
\usepackage{eurosym}
\usetikzlibrary {datavisualization}

\begin{document}

\pagestyle{fancy}
\fancyhead[L]{Terminale STMG}
\fancyhead[C]{\textbf{TD n°5 : statistiques à deux variables}}
\fancyhead[R]{\today}

\exe{}{
Une entreprise décide de mettre un nouveau produit sur le marché. Pour déterminer son prix de vente, elle réalise plusieurs tours
de vente limités en faisant varier le prix unitaire. Les chiffres de vente sont donnés dans le tableau ci-dessous :
 \begin{center}
  \begin{tabular}{|l|*5{c|}}
	 \hline
	 \textbf{Prix : $x_i$} & 50 & 55 & 60 & 80 & 100 \\
	 \hline
	 \textbf{Nombre ventes : $y_i$} & 150 & 138 & 130 & 92 & 48 \\
	 \hline
	\end{tabular}
 \end{center}
 
On a représenté dans le repère ci-dessous le nuage de points associé à cette série statistique à deux variables :
  
\begin{tikzpicture}
\begin{axis}%
    [
        grid=major,  
        x=2mm,
        y=1mm,
        xtick={50,55,...,110},   
        xmin=50,
        xmax=110,
        xlabel={\small $x$},
        axis x line=middle,
        ytick={40,50,...,160},
        tick label style={font=\small},
        ymin=40,
        ymax=160,
        ylabel={\small $y$},
        axis y line=middle,
        no markers,
        samples=100,
    ]
\addplot[teal, only marks] 
table[x=x,y=y, col sep=comma] {data_TD5.csv};
\end{axis}
\end{tikzpicture}

\begin{enumerate}
  \item Calculer les coordonnées du point moyen G de ce nuage de points. Placer G sur le graphique.
  \item \begin{enumerate}
	 \item Justifier qu'il est pertinent de réaliser un ajustement affine pour ce nuage de points.
	 \item Déterminer une équation de la droite $\Delta$ d'ajustement affine de $y$ en $x$ obtenue par la méthode des moindres carrés, sous la forme $y=ax+b$.
	 On arrondira les paramètres à $10^{-1}$ près.
	\end{enumerate}
	\item Vérifier que la droite $\Delta$ passe par G.
	\item À l'aide de cet ajustement affine, estimer le nombre d'offres vendues, si le prix unitaire est fixé à 70 euros.
	\item Selon vous, quel est le prix permettant à l'entreprise de réaliser le plus grand chiffre d'affaire ?
 \end{enumerate}
 
}{exe:1}{

}

\newpage

\exe{}{
La coopérative de Mai (à Clermont) décide de mesurer le volume sonore lors des concerts qu'elle organise sur un trimestre.
Pour cela, deux variables sont mesurées :
\begin{enumerate}[label=(\roman*)]
\item la pression accoustique (exprimée en pascals : Pa) ;
\item le niveau d'intensité sonore du bruit (exprimée en décibels : dB).
\end{enumerate}
Les résultats obtenus sont présentés dans le tableau ci-dessous :
 \begin{center}
  \begin{tabular}{|l|*8{c|}}
	 \hline
	 \textbf{Pression : $p_i$} & 0.5 & 1 & 3 & 5 & 7 & 10 & 13 & 15 \\
	 \hline
	 \textbf{Intensité : $y_i$} & 88 & 94 & 103 & 108 & 111 & 114 & 116 & 117 \\
	 \hline
	\end{tabular}
 \end{center}
 
Le nuage de points correspondant est donné ci-dessous :
\begin{center}
 \begin{tikzpicture}
\begin{axis}%
    [
        grid=major,  
        x=5mm,
        y=2mm,
        xtick={0,1,...,16},   
        xmin=0,
        xmax=20,
        xlabel={\small $p$},
        axis x line=middle,
        ytick={80,85,...,120},
        tick label style={font=\small},
        ymin=80,
        ymax=120,
        ylabel={\small $y$},
        axis y line=middle,
        no markers,
        samples=100,
    ]
\addplot[teal, only marks] 
table[x=x,y=y, col sep=comma] {data_TD5_2.csv};
\end{axis}
\end{tikzpicture}
\end{center}

\begin{enumerate}
\item Un ajustement affine du nuage de points vous semble-t-il pertinent ? Justifier.
\item On propose de réaliser une transformation logarithmique de la variable $p$ de pression en posant $x_i = \log(p_i)$.
Compléter le tableau ci-dessous en arrondissant $x_i$ à $10^{-2}$ près.
 \begin{center}
  \begin{tabular}{|l|*8{c|}}
	 \hline
	 \textbf{$x_i = \log(p_i)$} & \dots & \dots & \dots & \dots & \dots & 1 & \dots & \dots \\
	 \hline
	 \textbf{Intensité : $y_i$} & 88 & 94 & 103 & 108 & 111 & 114 & 116 & 117 \\
	 \hline
	\end{tabular}
 \end{center}
\item Représenter le nouveau nuage de points $M_i(x_i,y_i)$ dans un repère orthogonal.
\item Calculer les coordonnées du point moyen G de ce nuage de points et placer le sur le graphique.
\item Déterminer une équation de la droite $\Delta$ d'ajustement affine de $y$ en $x$ obtenue par la méthode des moindres carrés,
sous la forme $y=ax+b$. On arrondira les paramètres à $10^{-2}$ près. Tracer cette droite dans le repère.
\item Lors d'un concert de LuLu. (artiste Clermontois) l'oreille d'un spectateur est soumise à la pression de 130 Pa. Déterminer l'intensité
sonore atteinte lors de ce concert (au décibel près).
\end{enumerate}
}{exe:2}{

}

%%%%%%%%%%%%%%%%%%%%%%%%%%%%%

%\newpage
%\fancyhead[C]{\textbf{Solutions}}
%\shipoutAnswer

\end{document}